%%%%%%%%%%  scan idx 532  |  folio 519  %%%%%%%%%%
\SGAPrintPage{519}%

\sgaStableTarget{sga6.e10.expose}\section*{\underline{EXPOSÉ X}\\[3pt] \underline{\underline{FORMALISM OF INTERSECTIONS ON PROPER ALGEBRAIC SCHEMES}}}
\begin{center}
\underline{by O. Jussila}\\
(with an appendix by A. Grothendieck)
\end{center}

\sgaStableTarget{sga6.e10.sec.1}\section*{1. Compatibility of the filtrations with the composition law}
\[
K^{\bullet}(X)\times K_\bullet(X)\longrightarrow K_\bullet(X).
\]

\sgaStableTarget{sga6.e10.sec.1.1}\subsection*{1.1.}
First consider a noetherian scheme $X$ and the group $K_\bullet(X)$ of classes of coherent $\underline{O}_X$-modules. For every coherent $\underline{O}_X$-module $\underline{F}$, we shall denote its class in $K_\bullet(X)$ by $\cl_\bullet(\underline{F})$.

\sgaStableTarget{sga6.e10.def.1-1-1}\textbf{Definition 1.1.1.} \uline{For every integer} $n\in\mathbb Z$, \uline{let} $\Filt_n(X)$ \uline{denote the subgroup of} $K_\bullet(X)$ \uline{generated by the classes} $\cl_\bullet(\underline{F})$ \uline{of coherent} $\underline{O}_X$\uline{-modules} $\underline{F}$ \uline{such that}
\[
\dim(\Supp(\underline{F}))\leq n.
\]
This defines an increasing filtration $(\Filt_n(X))_{n\in\mathbb Z}$ of $K_\bullet(X)$ such that $\Filt_n(X)=0$ for $n<0$.

\sgaStableTarget{sga6.e10.prop.1-1-2}\textbf{Proposition 1.1.2.} \uline{Let} $X$ \uline{be a noetherian scheme,} $\underline{F}$ \uline{a coherent} $\underline{O}_X$\uline{-module such that} $\dim(\Supp(\underline{F}))\leq n$\uline{, and} $X(n)$ ($n\in\mathbb Z$\uline{) the subset} $\{x\in X\,;\,\dim(\overline{\{x\}})=n\}$ \uline{of} $X$\uline{. For every} $x\in X$, \uline{let} $Z(x)$ \uline{denote the integral closed subscheme associated with} $\overline{\{x\}}$\uline{. Then the length of the} $\underline{O}_{X,x}$\uline{-module} $\underline{F}_x$ \uline{is finite for every} $x\in X(n)$ \uline{and nonzero for only finitely many of the} $x\in X(n)$\uline{, and one has}


%%%%%%%%%%  scan idx 533  |  folio 520  %%%%%%%%%%
\SGAPrintPage{520}%

\sgaStableTarget{sga6.e10.eq.1-1-3}\begin{equation}\tag{1.1.3}
\cl_\bullet(\underline{F})=\sum_{x\in X(n)}\longg(\underline{F}_x)\cl_\bullet(\underline{O}_{Z(x)})+w
\end{equation}
\uline{where} $w\in\Filt_{n-1}(X)$.

By immediate induction on $n$, one obtains:

\sgaStableTarget{sga6.e10.cor.1-1-4}\textbf{Corollary 1.1.4.} \uline{For every noetherian scheme} $X$ \uline{and every integer} $n$, $\Filt_n(X)$ \uline{is the subgroup of} $K_\bullet(X)$ \uline{generated by the classes} $\cl_\bullet(\underline{O}_Z)$\uline{, where} $Z$ \uline{runs through the set of integral closed subschemes of dimension} $\leq n$ \uline{in} $X$.

\emph{Proof of} 1.1.2: Let $X'=\Supp(\underline{F})$. Since $X'$ is a noetherian space of dimension $\leq n$, the elements of $X'\cap X(n)$ are maximal in $X'$, hence finite in number, and $\longg(\underline{F}_x)$ is finite for every $x\in X'\cap X(n)$. This proves the first assertion.

Let ${K}=\Coh(X)$ be the category of coherent $\underline{O}_X$-modules, and let ${K}'$ be the subset of $\operatorname{Ob}({K})$ consisting of the coherent $\underline{O}_X$-modules concentrated on $X'$ that satisfy condition (1.1.3). Clearly ${K}'$ contains the sheaves $\underline{O}_{Z(x)}$ with $x\in X'$. Thus it is enough to show that ${K}'$ is an ``exact'' subset of $\operatorname{Ob}({K})$, since the dévissage lemma (EGA III 3.1.2) then gives $\underline{F}\in{K}'$.

Thus let
\sgaStableTarget{sga6.e10.eq.1-1-5}\begin{equation}\tag{1.1.5}
0\longrightarrow \underline{G}'\longrightarrow G\longrightarrow \underline{G}''\longrightarrow0
\end{equation}
be an exact sequence of $\underline{O}_X$-modules with two of its terms in ${K}'$. Then the sheaves $\underline{G}'$, $\underline{G}$ and $\underline{G}''$ are all concentrated on $X'$, so the lengths $\ell'_x$, $\ell_x$ and $\ell''_x$ of the $\underline{O}_{X,x}$-modules $\underline{G}'_x$, $\underline{G}_x$ and $\underline{G}''_x$ are finite


%%%%%%%%%%  scan idx 534  |  folio 521  %%%%%%%%%%
\SGAPrintPage{521}%

for every $x\in X(n)$, and one has
\sgaStableTarget{sga6.e10.eq.1-1-6}\begin{equation}\tag{1.1.6}
\ell_x=\ell'_x+\ell''_x\qquad (x\in X(n)).
\end{equation}
Assertion (1.1.3) for \emph{all} the terms of (1.1.5) then follows immediately from the hypothesis and the equation
\sgaStableTarget{sga6.e10.eq.1-1-7}\begin{equation}\tag{1.1.7}
\cl_\bullet(\underline{G})=\cl_\bullet(\underline{G}')+\cl_\bullet(\underline{G}'').
\end{equation}

It is known that $K_\bullet$ is a covariant functor on the category of locally noetherian schemes with respect to proper morphisms. The following proposition asserts that our filtration is functorial:

\sgaStableTarget{sga6.e10.prop.1-1-8}\textbf{Proposition 1.1.8.} \uline{Let} $f:X'\to X$ \uline{be a proper morphism of locally noetherian schemes, and let}
\[
f_*:K_\bullet(X')\longrightarrow K_\bullet(X)
\]
\uline{be the morphism induced by} $f$\uline{. Then}
\sgaStableTarget{sga6.e10.eq.1-1-9}\begin{equation}\tag{1.1.9}
f_*(\Filt_n(X'))\subset\Filt_n(X)
\end{equation}
\uline{for every integer} $n$.

\emph{Proof:} If $\underline{F}$ is a coherent $\underline{O}_{X'}$-module, then by definition (IV 2.11) one has:
\[
f_*(\cl_\bullet(\underline{F}))=\sum_{i\in\mathbb N}(-1)^i\cl_\bullet(R^if_*(\underline{F})).
\]
Now the $\underline{O}_X$-modules $R^if_*(\underline{F})$ are concentrated on $f(\Supp(\underline{F}))$, and by (EGA IV 5.4.1 (ii)) one has
\[
\dim(\Supp(\underline{F}))\geq\dim(f(\Supp(\underline{F}))),
\]
whence the assertion.


%%%%%%%%%%  scan idx 535  |  folio 522  %%%%%%%%%%
\SGAPrintPage{522}%

\sgaStableTarget{sga6.e10.sec.1.2}\subsection*{1.2.}
Now let $X$ be an arbitrary scheme, and consider the group $K^{\bullet}(X)$ of classes of locally free $\underline{O}_X$-modules (henceforth understood to be of finite type). It is an augmented pre-$\lambda$-algebra over the ring $H^0(X,\mathbb Z)$ (V 3.9.1). The map $f\rightsquigarrow f\cdot1\in K^{\bullet}(X)$ allows us henceforth to identify $H^0(X,\mathbb Z)$ with a subring of $K^{\bullet}(X)$. If $\underline{E}$ is a locally free $\underline{O}_X$-module, its class in $K^{\bullet}(X)$ will be denoted by $\cl^{\bullet}(\underline{E})$. Define the operators $\lambda^k$ ($k\in\mathbb N$) by
\sgaStableTarget{sga6.e10.eq.1-2-1}\begin{equation}\tag{1.2.1}
\lambda^k(\cl^{\bullet}(\underline{E}))=\cl^{\bullet}(\Lambda_{\underline{O}}^k\underline{E})
\end{equation}
and by
\sgaStableTarget{sga6.e10.eq.1-2-2}\begin{equation}\tag{1.2.2}
\lambda^k(x+y)=\sum_{\substack{i+j=k\\ i,j\in\mathbb N}}\lambda^i(x)\lambda^j(y).
\end{equation}
(V 2.2). Recall that in the $\lambda$-filtration $(\Filt^i(X))_{i\in\mathbb Z}$ of $K^{\bullet}(X)$, one has $\Filt^i(X)=K^{\bullet}(X)$ for $i\leq0$, and that for $i\geq1$, $\Filt^i(X)$ is the ideal generated by the products
\sgaStableTarget{sga6.e10.eq.1-2-3}\begin{equation}\tag{1.2.3}
P=\gamma^{i_1}(x_1)\cdots\gamma^{i_k}(x_k)
\end{equation}
with $1\leq i_1,\ldots,i_k$; $i_1+\cdots+i_k\geq i$ and $x_1,\ldots,x_k\in I=\Ker(\varepsilon)$ ($\varepsilon$ being the augmentation homomorphism, cf. Exposé V 3.10). It was proved (V 3.10) that $\Filt^i(X)$ is even the $H^0(X,\mathbb Z)$-submodule (hence the abelian subgroup, when $X$ is connected) generated by the products (1.2.3), where one may moreover suppose that the elements $x_p$ ($p=1,\ldots,k$) have the form
\sgaStableTarget{sga6.e10.eq.1-2-4}\begin{equation}\tag{1.2.4}
x_p=\cl^{\bullet}(\underline{E}_p)-\underline{n}_p,
\end{equation}
with $\underline{E}_p$ a locally free $\underline{O}_X$-module of rank $\underline{n}_p\in H^0(X,\mathbb Z)$.


%%%%%%%%%%  scan idx 536  |  folio 523  %%%%%%%%%%
\SGAPrintPage{523}%

\sgaStableTarget{sga6.e10.sec.1.3}\subsection*{1.3.}
Recall (IV 2.10) that the bifunctor
\[
(\underline{E},\underline{F})\rightsquigarrow \underline{E}\otimes_{\underline{O}_X}\underline{F},
\]
where $\underline{E}$ is a locally free $\underline{O}_X$-module and $\underline{F}$ a coherent $\underline{O}_X$-module, induces a composition law
\sgaStableTarget{sga6.e10.eq.1-3-1}\begin{equation}\tag{1.3.1}
K^{\bullet}(X)\times K_\bullet(X)\longrightarrow K_\bullet(X)
\end{equation}
that makes $K_\bullet(X)$ a unital $K^{\bullet}(X)$-module. The key result of the present exposé is the following

\sgaStableTarget{sga6.e10.thm.1-3-2}\textbf{Theorem 1.3.2.} \uline{Let} $X$ \uline{be a noetherian scheme. For all} $i,j\in\mathbb Z$, \uline{the composition law} (1.3.1) \uline{induces by restriction a composition law}
\sgaStableTarget{sga6.e10.eq.1-3-3}\begin{equation}\tag{1.3.3}
\Filt^i(X)\times\Filt_j(X)\longrightarrow\Filt_{j-i}(X).
\end{equation}

\sgaStableTarget{sga6.e10.cor.1-3-3}\textbf{Corollary 1.3.3.} \uline{If} $X$ \uline{is a noetherian scheme of dimension} $\leq j$\uline{, then the canonical homomorphism}
\[
\Theta_X:K^{\bullet}(X)\longrightarrow K_\bullet(X)
\]
\uline{defined by} $\Theta_X(x)=x\cdot\cl_\bullet(\underline{O}_X)$ ($x\in K^{\bullet}(X)$\uline{) induces for every} $i\in\mathbb Z$ \uline{a homomorphism}
\[
\Filt^i(X)\longrightarrow\Filt_{j-i}(X).
\]

\emph{Proof:} First let us show that (1.3.3), for $i+j\leq k\in\mathbb Z$ and for \emph{every} integral closed subscheme of dimension $\leq j$, implies (1.3.2) for $i+j\leq k$. By (1.1.4), it is enough to show that for every $x\in\Filt^i(X)$ and for \emph{every} integral closed subscheme $Z$ of dimension $\leq j$ one has

